\documentclass[12pt,a4paper]{article}
\usepackage[T2A]{fontenc}
\usepackage[utf8]{inputenc}
\usepackage[russian]{babel}
\usepackage{geometry}
\usepackage{tabularx} 
\geometry{margin=2.5cm}
\usepackage{amsmath,amssymb,mathtools}
\usepackage{enumitem}
\usepackage{tabularx}
\usepackage{hyperref}
\hypersetup{colorlinks=true,linkcolor=black,urlcolor=blue}
\setlength{\parindent}{0pt}
\setlength{\parskip}{6pt}
\begin{document}
{\Large\bfseries
\begin{flushleft}
Отчет по лабораторной работе №3 (Вариант №21)
\end{flushleft}}
\vspace{0.5em}\hrule\vspace{1.2em}

\section*{Задание}
\begin{itemize}
  \item Проанализировать язык на детермизим. Построить PDA: в случае детерминированного языка построить DPDA, в случае недетерминированного языка - PDA без избыточного недетерминизма (для каждого недетерминированного перехода привести пример слов с общим длинным префиксом, в которых на переходе существенно отличается поведение стека). 
  \item Проанализировать язык на беспрефиксность.
  \item В случае детерминированного языка - проанализировать его на LL-свойство (дополнительное задание).
  \item Построить КС-грамматику для языка, если таковой нет (в случае LL-языка - удобно использовать LL(k)-грамматику). Пересечь её с регулярными аппроксимациями сверху, полученными First, Follow множеств («LL(1)-автомат»), и позиционного автомата («LR(0)-автомат»).
\end{itemize}

\section*{Исходная грамматика G}
\[
\begin{aligned}
S &\to aTTa \\
T &\to SbbS \\
S &\to abba \\
T &\to bab \\
T &\to aTa
\end{aligned}
\]

\newpage

\section*{Детерминизм}
\subsection*{Доказательство недетерминированности}
Докажу, что данный язык недетерминированный. Рассмотрим пересечение данного языка с регулярным выражением
\[
R = a^{+}abba(bbabbababa)^{+}(bbabbaa^{+}baba|\epsilon)
\]

Возьму два слова, получив их следующим образом:

\[
\begin{aligned}
1)\quad \underline{S}
&\Rightarrow a\underline{T}Ta\\
&\Rightarrow aa^{n}\underline{T}a^{n}Ta\\
&\Rightarrow aa^{n}\underline{S}bbSa^{n}Ta\\
&\Rightarrow aa^{n}Sbb\underline{S}a^{n}\underline{T}a\\
&\Rightarrow aa^{n}\underline{S}bbabbaa^{n}baba\\
&\Rightarrow aa^{n}a\underline{T}Tabbabbaa^{n}baba\\
&\Rightarrow aa^{n}aSbb\underline{S}\underline{T}abbabbaa^{n}baba\\
&\Rightarrow aa^{n}(a)\underline{S}(bbabbababa)bbabbaa^{n}baba\\
&\Rightarrow aa^{n}(a)^{n}\underline{S}(bbabbababa)^{n}bbabbaa^{n}baba\\
&\Rightarrow aa^{n}a^{n}abba(bbabbababa)^{n}bbabbaa^{n}baba\\
&\Rightarrow a^{2n+1}abba(bbabbababa)^{n}bbabbaa^{n}baba.
\end{aligned}
\]

\[
\begin{aligned}
2)\quad \underline{S}
&\Rightarrow a\underline{T}Ta\\
&\Rightarrow aSbb\underline{ST}a\\
&\Rightarrow (a)\underline{S}(bbabbababa)\\
&\Rightarrow a^{2n+1}\underline{S}(bbabbababa)^{2n+1}\\
&\Rightarrow a^{2n+1}abba(bbabbababa)^{2n+1}\\
&\Rightarrow a^{2n+1}abba(bbabbababa)^{n}(bbabbababa)^{n+1}.
\end{aligned}
\]

Оба слова имеют общий префикс $a^{2n+1}abba(bbabbababa)^{n}$. Таким образом получается:
\[
\begin{array}{l|l}
a^{2n+1}abba(bbabbababa)^{n} & bbabbaa^{n}baba\\
a^{2n+1}abba(bbabbababa)^{n} & (bbabbababa)^{n+1}
\end{array}
\]

1) Для первого случая накачиваем \((bbabbababa)^n\) и \(a^n\). После накачки получается \(a^{2n+1}abba(bbabbababa)^{n+i}bbabbaa^{n+k}baba\). В резульате теряется равенство с блоком \(a^{2n+1}\), и слово выводится из языка. 

2) Для второго случая накачиваем \((bbabbababa)^n\) и \((bbabbababa)^{n+1}\). После накачки получается \(a^{2n+1}abba(bbabbababa)^{2n+1+i}\) В резульате теряется равенство с блоком \(a^{2n+1}\), и слово выводится из языка.

Следовательно, \(L\cap R\) не является DCFL. Тогда и сам язык \(L\) не является DCFL. А значит для \(L\) нельзя построить DPDA

\subsection*{PDA}
Построю PDA, распознающий грамматику G:
\begin{center}
\begin{minipage}{\linewidth}
\centering
\includegraphics[width=1.0\linewidth]{pda.jpg}
\captionof{Рис. 1: Стековый автомат (PDA)}
\label{fig:pda.jpg}
\end{minipage}
\end{center}

\subsection*{Недетерминированные переходы и примеры слов, демонстрирующих отличия в поведении стека}
1) Недетерминированные переходы из состояния 1:
\begin{itemize}
\item переход по $a$ с операцией $T/TA$ в состояние $1$
\item переход по $a$ с операцией $T/BBS$ в состояние $2$
\end{itemize}

\[
\begin{array}{l|l|l}
aa^{2n} & baba^{2n}baba & \text{(для первого указанного перехода)}\\
aa^{2n} & (ababaaabbabb)^{n}ababa & \text{(для второго указанного перехода)}
\end{array}
\]

2) Недетерминированные переходы из состояния 2:
\begin{itemize}
\item переход по $a$ c операцией $x/TATAx$ в состояние $1$
\item переход по $a$ с операцией $x/BBSTAx$ в состояние $2$
\end{itemize}

\[
\begin{array}{l|l|l}
aa^{2n}bba & (bbabbaababa)^{n} & \text{(для первого указанного перехода)}\\
aa^{2n}bba & (bbabbababa)^{2n} & \text{(для второго указанного перехода)}
\end{array}
\]

\section*{Беспрефиксность}
Рассмотрим слово $w_1=aababababa$. Слово $w_1$ принадлежит языку и выводится следующим образом:
\[
\begin{aligned}
S &\Rightarrow a\underline{T}Ta \\
  &\Rightarrow aa\underline{T}aTa \\
  &\Rightarrow aababa\underline{T}a \\
  &\Rightarrow aababababa
\end{aligned}
\]
При добавлении $abbabbababa$ в виде суффикса к $w_1$ получится $w_2=aababababaabbabbababa$.
Слово $w_2$ также принадлежит языку и выводится следующим образом:
\[
\begin{aligned}
S &\Rightarrow a\underline{T}Ta \\
  &\Rightarrow a\underline{S}bbSTa \\
  &\Rightarrow aa\underline{T}TabbSTa \\
  &\Rightarrow aabab\underline{T}abbSTa \\
  &\Rightarrow aababa\underline{T}aabbSTa \\
  &\Rightarrow aababababaabb\underline{S}Ta \\
  &\Rightarrow aababababaabbabba\underline{T}a \\
  &\Rightarrow aababababaabbabbababa
\end{aligned}
\]

Таким образом, $w_1$ является префиксом $w_2$. Значит, язык, задаваемый этой КС-грамматикой, не является беспрефиксным.

\section*{LL-свойства}
LL-свойство требует, чтобы язык был детерминированным.\\
Мой язык недетерминирован, следовательно этот язык не обладает LL-свойством

\section*{Построение аппроксимаций}
\subsection*{LL(1)-аппроксимация}
Для моей грамматики найду First, Follow, Last множества:
\[
\mathrm{First}(L(G))=\{a_{1},\,a_{3}\}
\]

\[
\mathrm{Follow}(L(G))=\left\{
\begin{aligned}
& a_{1}a_{7},\,a_{1}a_{1},\,a_{1}a_{3},\,a_{1}b_{11},\\
& a_{2}a_{7},\,a_{2}a_{1},\,a_{2}a_{3},\,a_{2}b_{11},\,a_{2}a_{2},\,a_{2}b_{9},\,a_{2}a_{8},\\
& a_{3}b_{4},\\
& b_{4}b_{5},\\
& b_{5}a_{6},\\
& a_{6}a_{7},\,a_{6}a_{1},\,a_{6}a_{3},\,a_{6}b_{11},\,a_{6}a_{2},\,a_{6}b_{9},\,a_{6}a_{8},\\
& a_{7}a_{7},\,a_{7}a_{1},\,a_{7}a_{3},\,a_{7}b_{11},\\
& a_{8}a_{7},\,a_{8}a_{1},\,a_{8}a_{3},\,a_{8}b_{11},\,a_{8}a_{2},\,a_{8}a_{8},\\
& b_{9}b_{10},\\
& b_{10}a_{1},\,b_{10}a_{3},\\
& b_{11}a_{12},\\
& a_{12}b_{13},\\
& b_{13}a_{7},\,b_{13}a_{1},\,b_{13}a_{3},\,b_{13}b_{11},\,b_{13}a_{2},\,b_{13}a_{8}
\end{aligned}
\right\}
\]

\[
\mathrm{Last}(L(G))=\{a_{2},\,a_{6}\}
\]

На основе вычисленных множеств получаю следующий НКА:
\begin{center}
\begin{minipage}{\linewidth}
\centering
\includegraphics[width=1.0\linewidth]{ll1nka.jpg}
\captionof{Рис. 2: LL(1)-аппроксимация сверху}
\label{fig:ll1nka.jpg}
\end{minipage}
\end{center}

\begin{center}
\begin{minipage}{\linewidth}
\centering
\includegraphics[width=1.0\linewidth]{ll1nkamin.jpg}
\captionof{Рис. 3: LL(1)-аппроксимация сверху (минимизированная)}
\label{fig:ll1nkamin.jpg}
\end{minipage}
\end{center}

\subsection*{Правила из пересечения грамматики и LL(1)-аппроксимации}
\[
\begin{aligned}
\langle1,T,1\rangle&\to a\langle1,T,1\rangle a\\
\langle1,T,1\rangle&\to a\langle1,T,8\rangle a\\
\langle3,T,1\rangle&\to a\langle3,T,1\rangle a\\
\langle3,T,1\rangle&\to a\langle3,T,8\rangle a\\
\langle3,T,1\rangle&\to a\langle3,T,10\rangle a\\
\langle8,T,1\rangle&\to a\langle1,T,1\rangle a\\
\langle8,T,1\rangle&\to a\langle1,T,8\rangle a\\
\langle10,T,1\rangle&\to a\langle1,T,1\rangle a\\
\langle10,T,1\rangle&\to a\langle1,T,8\rangle a\\
\langle1,T,1\rangle&\to \langle1,S,1\rangle bb\langle7,S,1\rangle\\
\langle3,T,1\rangle&\to \langle3,S,1\rangle bb\langle7,S,1\rangle\\
\langle3,T,1\rangle&\to \langle3,S,10\rangle bb\langle0,S,1\rangle\\
\langle3,T,10\rangle&\to \langle3,S,10\rangle bb\langle0,S,10\rangle\\
\langle8,T,1\rangle&\to \langle8,S,1\rangle bb\langle7,S,1\rangle\\
\langle10,T,1\rangle&\to \langle10,S,1\rangle bb\langle7,S,1\rangle\\
\langle1,T,8\rangle&\to bab\\
\langle3,T,8\rangle&\to bab\\
\langle8,T,8\rangle&\to bab\\
\langle10,T,8\rangle&\to bab\\
\langle0,S,1\rangle&\to a\langle3,T,1\rangle\langle1,T,1\rangle a\\
\langle0,S,1\rangle&\to a\langle3,T,1\rangle\langle1,T,8\rangle a\\
\langle0,S,1\rangle&\to a\langle3,T,8\rangle\langle8,T,1\rangle a\\
\langle0,S,1\rangle&\to a\langle3,T,8\rangle\langle8,T,8\rangle a\\
\langle0,S,1\rangle&\to a\langle3,T,10\rangle\langle10,T,1\rangle a\\
\langle0,S,1\rangle&\to a\langle3,T,10\rangle\langle10,T,8\rangle a\\
\langle1,S,1\rangle&\to a\langle1,T,1\rangle\langle1,T,1\rangle a\\
\langle1,S,1\rangle&\to a\langle1,T,1\rangle\langle1,T,8\rangle a\\
\langle1,S,1\rangle&\to a\langle1,T,8\rangle\langle8,T,1\rangle a\\
\langle1,S,1\rangle&\to a\langle1,T,8\rangle\langle8,T,8\rangle a\\
\end{aligned}
\]

\[
\begin{aligned}
\langle3,S,1\rangle&\to a\langle3,T,1\rangle\langle1,T,1\rangle a\\
\langle3,S,1\rangle&\to a\langle3,T,1\rangle\langle1,T,8\rangle a\\
\langle3,S,1\rangle&\to a\langle3,T,8\rangle\langle8,T,1\rangle a\\
\langle3,S,1\rangle&\to a\langle3,T,8\rangle\langle8,T,8\rangle a\\
\langle3,S,1\rangle&\to a\langle3,T,10\rangle\langle10,T,1\rangle a\\
\langle3,S,1\rangle&\to a\langle3,T,10\rangle\langle10,T,8\rangle a\\
\langle7,S,1\rangle&\to a\langle1,T,1\rangle\langle1,T,1\rangle a\\
\langle7,S,1\rangle&\to a\langle1,T,1\rangle\langle1,T,8\rangle a\\
\langle7,S,1\rangle&\to a\langle1,T,8\rangle\langle8,T,1\rangle a\\
\langle7,S,1\rangle&\to a\langle1,T,8\rangle\langle8,T,8\rangle a\\
\langle8,S,1\rangle&\to a\langle1,T,1\rangle\langle1,T,1\rangle a\\
\langle8,S,1\rangle&\to a\langle1,T,1\rangle\langle1,T,8\rangle a\\
\langle8,S,1\rangle&\to a\langle1,T,8\rangle\langle8,T,1\rangle a\\
\langle8,S,1\rangle&\to a\langle1,T,8\rangle\langle8,T,8\rangle a\\
\langle10,S,1\rangle&\to a\langle1,T,1\rangle\langle1,T,1\rangle a\\
\langle10,S,1\rangle&\to a\langle1,T,1\rangle\langle1,T,8\rangle a\\
\langle10,S,1\rangle&\to a\langle1,T,8\rangle\langle8,T,1\rangle a\\
\langle10,S,1\rangle&\to a\langle1,T,8\rangle\langle8,T,8\rangle a\\
\langle0,S,10\rangle&\to abba\\
\langle1,S,1\rangle&\to abba\\
\langle3,S,10\rangle&\to abba\\
\langle7,S,1\rangle&\to abba\\
\langle8,S,1\rangle&\to abba\\
\langle10,S,1\rangle&\to abba
\end{aligned}
\]

\newpage

\subsection*{LR(0)-аппроксимация}
\begin{center}
\begin{minipage}{\linewidth}
\centering
\includegraphics[width=1.0\linewidth]{lr0.jpg}
\captionof{Рис. 4: Позиционный автомат}
\label{fig:lr0.jpg}
\end{minipage}
\end{center}

\begin{center}
\begin{minipage}{\linewidth}
\centering
\includegraphics[width=1.0\linewidth]{promezhpredst.jpg}
\captionof{Рис. 5: Промежуточное представление позиционного автомата}
\label{fig:promezhpredst.jpg}
\end{minipage}
\end{center}

\begin{center}
\begin{minipage}{\linewidth}
\centering
\includegraphics[width=1.0\linewidth]{epspereh.jpg}
\captionof{Рис. 6: Замена переходов по нетерминалам на $\epsilon$-переходы}
\label{fig:epspereh.jpg}
\end{minipage}
\end{center}

\begin{center}
\begin{minipage}{\linewidth}
\centering
\includegraphics[width=1.0\linewidth]{lr0min.jpg}
\captionof{Рис. 7: LR(0)-аппроксимация сверху}
\label{fig:lr0min.jpg}
\end{minipage}
\end{center}

\subsection*{Правила из пересечения грамматики и LR(0)-аппроксимации}
\[
\begin{aligned}
\langle1,T,1\rangle&\to a\langle1,T,1\rangle a\\
\langle1,T,1\rangle&\to a\langle1,T,7\rangle a\\
\langle7,T,1\rangle&\to a\langle1,T,1\rangle a\\
\langle7,T,1\rangle&\to a\langle1,T,7\rangle a\\
\langle8,T,1\rangle&\to a\langle1,T,1\rangle a\\
\langle8,T,1\rangle&\to a\langle1,T,7\rangle a\\
\langle10,T,1\rangle&\to a\langle10,T,1\rangle a\\
\langle10,T,1\rangle&\to a\langle10,T,7\rangle a\\
\langle10,T,1\rangle&\to a\langle10,T,8\rangle a\\
\langle1,T,1\rangle&\to \langle1,S,1\rangle bb\langle6,S,1\rangle\\
\langle7,T,1\rangle&\to \langle7,S,1\rangle bb\langle6,S,1\rangle\\
\langle8,T,1\rangle&\to \langle8,S,1\rangle bb\langle6,S,1\rangle\\
\langle10,T,1\rangle&\to \langle10,S,1\rangle bb\langle6,S,1\rangle\\
\end{aligned}
\]

\[
\begin{aligned}
\langle10,T,1\rangle&\to \langle10,S,8\rangle bb\langle0,S,1\rangle\\
\langle10,T,8\rangle&\to \langle10,S,8\rangle bb\langle0,S,8\rangle\\
\langle1,T,7\rangle&\to bab\\
\langle7,T,7\rangle&\to bab\\
\langle8,T,7\rangle&\to bab\\
\langle10,T,7\rangle&\to bab\\
\langle0,S,1\rangle&\to a\langle10,T,1\rangle\langle1,T,1\rangle a\\
\langle0,S,1\rangle&\to a\langle10,T,1\rangle\langle1,T,7\rangle a\\
\langle0,S,1\rangle&\to a\langle10,T,7\rangle\langle7,T,1\rangle a\\
\langle0,S,1\rangle&\to a\langle10,T,7\rangle\langle7,T,7\rangle a\\
\langle0,S,1\rangle&\to a\langle10,T,8\rangle\langle8,T,1\rangle a\\
\langle0,S,1\rangle&\to a\langle10,T,8\rangle\langle8,T,7\rangle a\\
\langle1,S,1\rangle&\to a\langle1,T,1\rangle\langle1,T,1\rangle a\\
\langle1,S,1\rangle&\to a\langle1,T,1\rangle\langle1,T,7\rangle a\\
\langle1,S,1\rangle&\to a\langle1,T,7\rangle\langle7,T,1\rangle a\\
\langle1,S,1\rangle&\to a\langle1,T,7\rangle\langle7,T,7\rangle a\\
\langle6,S,1\rangle&\to a\langle1,T,1\rangle\langle1,T,1\rangle a\\
\langle6,S,1\rangle&\to a\langle1,T,1\rangle\langle1,T,7\rangle a\\
\langle6,S,1\rangle&\to a\langle1,T,7\rangle\langle7,T,1\rangle a\\
\langle6,S,1\rangle&\to a\langle1,T,7\rangle\langle7,T,7\rangle a\\
\langle7,S,1\rangle&\to a\langle1,T,1\rangle\langle1,T,1\rangle a\\
\langle7,S,1\rangle&\to a\langle1,T,1\rangle\langle1,T,7\rangle a\\
\langle7,S,1\rangle&\to a\langle1,T,7\rangle\langle7,T,1\rangle a\\
\langle7,S,1\rangle&\to a\langle1,T,7\rangle\langle7,T,7\rangle a\\
\langle8,S,1\rangle&\to a\langle1,T,1\rangle\langle1,T,1\rangle a\\
\langle8,S,1\rangle&\to a\langle1,T,1\rangle\langle1,T,7\rangle a\\
\langle8,S,1\rangle&\to a\langle1,T,7\rangle\langle7,T,1\rangle a\\
\langle8,S,1\rangle&\to a\langle1,T,7\rangle\langle7,T,7\rangle a\\
\langle10,S,1\rangle&\to a\langle10,T,1\rangle\langle1,T,1\rangle a\\
\langle10,S,1\rangle&\to a\langle10,T,1\rangle\langle1,T,7\rangle a\\
\langle10,S,1\rangle&\to a\langle10,T,7\rangle\langle7,T,1\rangle a\\
\langle10,S,1\rangle&\to a\langle10,T,7\rangle\langle7,T,7\rangle a\\
\langle10,S,1\rangle&\to a\langle10,T,8\rangle\langle8,T,1\rangle a\\
\langle10,S,1\rangle&\to a\langle10,T,8\rangle\langle8,T,7\rangle a\\
\langle0,S,8\rangle&\to abba\\
\langle1,S,1\rangle&\to abba\\
\langle6,S,1\rangle&\to abba\\
\langle7,S,1\rangle&\to abba\\
\langle8,S,1\rangle&\to abba\\
\langle10,S,8\rangle&\to abba
\end{aligned}
\]

\end{document}